%%%%%%%%%%%%%%%%%%%%%%%%%%%%%%%%%%%%%%%%%
% Navid Safahieh -- Job Market CV (2026-27)
% Self-contained: standard article class, no custom .cls required.
%%%%%%%%%%%%%%%%%%%%%%%%%%%%%%%%%%%%%%%%%

\documentclass[11pt]{article}

\usepackage[left=0.75in,top=0.7in,right=0.75in,bottom=0.7in]{geometry}
\usepackage[T1]{fontenc}
\usepackage[utf8]{inputenc}
\usepackage{enumitem}
\usepackage{titlesec}
\usepackage{xcolor}
\usepackage{needspace}
\usepackage[hidelinks]{hyperref}

\hypersetup{
  colorlinks = true,
  urlcolor   = blue,
  linkcolor  = blue,
  citecolor  = blue
}

\pagestyle{empty}
\setlength{\parindent}{0pt}

% ---- Section headings: small caps with a rule underneath ----
\titleformat{\section}
  {\normalfont\large\scshape}
  {}{0em}{}[\vspace{-0.6em}\rule{\textwidth}{0.4pt}]
\titlespacing*{\section}{0pt}{0.85em}{0.45em}

% ---- Helper: entry with title flush left, date flush right ----
\newcommand{\entry}[2]{%
  \needspace{2\baselineskip}%
  \textbf{#1}\hfill #2\par}

% ---- Helper: role / institution line ----
\newcommand{\role}[2]{\textit{#1}\hfill #2\par\vspace{0.2em}}

% ---- Helper: paper title + description block ----
\newcommand{\paper}[1]{\needspace{3\baselineskip}\textbf{#1}\par\vspace{0.15em}}

% ---- Helper: current month and year, e.g. "September 2026" ----
\newcommand{\monthyear}{\ifcase\month\or January\or February\or March\or
  April\or May\or June\or July\or August\or September\or October\or
  November\or December\fi\ \number\year}

\newenvironment{desc}
  {\begin{itemize}[leftmargin=1.2em,topsep=0.1em,itemsep=0.15em,parsep=0pt]}
  {\end{itemize}\vspace{0.5em}}

\newenvironment{tight}
  {\begin{itemize}[leftmargin=1.2em,topsep=0.1em,itemsep=0.15em,parsep=0pt]}
  {\end{itemize}}

\begin{document}

% ======================= HEADER =======================
\begin{center}
  {\LARGE\bfseries Navid Safahieh}\\[0.5em]
  Department of Economics, Mitch Daniels School of Business\\
  Purdue University, 403 Mitch Daniels Blvd, West Lafayette, IN 47907\\[0.3em]
  +1 (765) 771--0021 \quad$\cdot$\quad
  \href{mailto:nsafahie@purdue.edu}{nsafahie@purdue.edu} \quad$\cdot$\quad
  \href{https://navidzky.github.io}{navidzky.github.io}\\[0.2em]
  Nationality: New Zealand / Iran\\[0.4em]
  \textbf{On the job market 2026--27}
\end{center}

% ======================= RESEARCH INTERESTS =======================
\section{Research Interests}
Transportation Economics, Search and Matching, Macroeconomics, Causal Inference

% ======================= EDUCATION =======================
\section{Education}

\entry{Ph.D. in Economics}{Expected 2027}
Purdue University, Mitch Daniels School of Business, West Lafayette, IN\par\vspace{0.4em}

\entry{M.Sc. in Economics}{2021}
University of Tehran, Tehran, Iran\par\vspace{0.4em}

\entry{B.Sc. in Civil and Environmental Engineering}{2018}
Sharif University of Technology, Tehran, Iran\par\vspace{0.5em}

\textbf{Dissertation Committee:} Cathy Zhang (Chair), Mario J. Crucini, Chong Xiang

% ======================= JOB MARKET PAPER =======================
\section{Job Market Paper}

\paper{Match Now, Arrive Later: Transport Markets with Remote Search}
In transport markets, carriers are often hired while still far from the cargo, and the customer
waits while the carrier travels toward the cargo. I develop a dynamic spatial directed search model
that allows for this remote matching. In existing models of transport markets, matching is local: a
carrier must sail to a cargo before it can compete for it, sinking the cost of the trip for a cargo it
may not win. With remote matching, a carrier competes from where it is for cargoes at every location,
and a customer can attract carriers from every location, at the cost of waiting while they travel, so
competition for carriers and for cargoes extends across locations. In a two-location example in which
cargo moves in one direction only, as trips become more valuable, carriers stop sailing to the cargo
before they are hired and wait to be hired remotely, even though being near the cargo remains more
valuable: the cost of the time a carrier spends sailing without a contract, instead of applying, grows
faster with the value of a trip than the premium of being near the cargo. Customers then hire more
distant carriers and wait longer, consistent with tankers being fixed farther from the loading port
when freight markets are strong. Compared with an otherwise identical model with local matching only,
once trips are valuable enough, carriers are worth more at every location and the premium of being
near the cargo is smaller. In an application with the Persian Gulf, Europe and Asia, I close Bab
el-Mandeb, which lengthens voyages from the Gulf to Europe. The outcome depends on the value of a trip
to Europe. When it is high enough that the detour is worth taking, the shock negatively affects the
route to Asia, which does not use the passage: its fee rises and less cargo moves on it. When it is
lower, carriers shift to the Asia route and its fee falls. Over a wide range of values, remote matching
amplifies the spillover: customers shipping to Europe keep shipping by hiring distant carriers, and
shipments to Asia fall more than with local matching.

% ======================= PUBLICATIONS =======================
\needspace{8\baselineskip}
\section{Publications}

\paper{\href{https://doi.org/10.1016/j.iatssr.2022.03.004}{Modeling the Accuracy of Traffic Crash Prediction Models}}
\begin{desc}
  \item \textit{IATSS Research}, 2022. With Mohammad Hesam Rashidi, Soheil Keshavarz, Parham Pazari, and Amir Samimi.
  \item Crash forecasting enables safety planners to take appropriate actions before casualty or loss
  occurs. Identifying and analyzing the attributes influencing forecasting accuracy is of great
  importance in road crash forecasting. This study aims to model the forecasting accuracy of 31
  provinces using their macroeconomic variables and road traffic indicators. Iran's road crashes
  throughout 2011--2018 are calibrated and cross-validated using the Holt-Winters (HW) forecasting
  method. The sensitivity of crash forecast reliability is studied by a regression model. The
  results suggested that the root mean square error (RMSE) of crash prediction increased among the
  provinces with higher and more variant average monthly crashes. On the contrary, the accuracy of
  crash prediction improved in provinces with higher per capita GDP, and higher traffic exposure. A
  1\% increase in crash variability, average historical crash count, GDP per capita, and traffic
  exposure, respectively, resulted in a 0.65\%, 0.52\%, $-$0.38\%, and $-$0.13\% change in the RMSE
  of forecasting. The addition of traffic exposure and macroeconomic factors significantly enhanced
  the model fit and improved the adjusted R-squared by 14\% compared to the reduced model that only
  used the historical average and variability of crash count as the independent variables. The
  findings of this research suggest planners and policymakers should consider the notable influence
  of macroeconomic factors and traffic indicators on the crash forecasting accuracy.
\end{desc}

% ======================= WORKING PAPERS =======================
\section{Working Papers}

\paper{\href{https://dx.doi.org/10.2139/ssrn.6379739}{When Outside Options Shape the Labor Market: A Directed Search Perspective}}
\begin{desc}
  \item This paper studies how differences in outside options across workers affect wages, firm entry, and
  job allocation in the labor market. I develop a directed search model with two types of
  workers---skilled and unskilled---and two types of firms---high-tech and low-tech. Only skilled
  workers have outside options, such as unemployment benefits or opportunities in other regions,
  creating an asymmetry between worker groups. The model extends Shi (2002) by allowing for outside
  options and by relaxing the assumption that high-tech firms hiring skilled workers are always more
  profitable. These changes generate several possible equilibrium outcomes that depend on the value
  of skilled workers' outside options and their share in the population. The results show that
  stronger outside options raise posted wages for skilled workers but can reduce high-tech firm
  entry. As fewer high-tech jobs are created, skilled workers' expected wages may fall, and
  unskilled workers are indirectly affected through changes in firm composition and wage
  competition. Overall, the paper shows that while better outside options make workers more
  selective, they can also discourage firm entry and lower overall wage opportunities in the market.
\end{desc}

\paper{\href{https://dx.doi.org/10.2139/ssrn.4830956}{The Heterogeneous Effects of Exogenous Shocks on Public Transportation Demand}}
\begin{desc}
  \item With Kowsar Yousefi and Amir Samimi.
  \item This study evaluates the impact of the gas price hike and the COVID-19 pandemic on subway usage in
  2019 in the Tehran and Alborz provinces. The hike was about 89\% and increased subway usage, while
  the COVID-19 pandemic reduced it. Access to exact hourly transactions over a year, we utilize an
  event study methodology to estimate the cross-price elasticity and the COVID-19 pandemic's effect.
  Results show that subway usage increased by about 3.9 percent due to raised gas prices, so the
  calculated cross-price elasticity will be about 0.044, and in hindsight, subway usage decreased by
  about 42.6 percent due to the COVID-19 pandemic. Furthermore, this research shows that the
  cross-price elasticity is higher when the share of the morning entries---commuted
  trips---increases, whereas the COVID-19 pandemic's impact is more potent as the share of
  noncommuted trips increases. Last but not least, the results show that the calculated cross-price
  elasticity tends to decrease, and the drop caused by COVID-19 tends to be more influential as we
  move toward wealthier regions. These results indicate that the impoverished groups are more prone
  to any change in prices, plus they cannot protect themselves against unpleasant circumstances such
  as a life-threatening pandemic.
\end{desc}

% ======================= PRESENTATIONS =======================
\section{Conference Presentations}
\begin{tight}
  \item MEA 90th Annual Meetings, Chicago, IL \hfill March 2026
\end{tight}

% ======================= TRAINING =======================
\section{Additional Training}
\begin{tight}
  \item Estimation of Structural Models with Conditional Choice Probabilities, Notre Dame, IN \hfill Summer 2025
  \item 14th Annual Research Design for Causal Inference, Northwestern University, IL \hfill Summer 2025
\end{tight}

% ======================= TEACHING =======================
\section{Teaching Experience}

\entry{ECON 252 Principles of Macroeconomics (Undergraduate)}{Fall 2025}
\role{Instructor}{Purdue University}

\entry{ECON 252 Principles of Macroeconomics (Undergraduate)}{Summer 2025}
\role{Instructor}{Purdue University}

\entry{Intermediate Macroeconomics (Undergraduate)}{Fall 2025}
\role{Teaching Assistant to Prof. Cathy Zhang}{Purdue University}

\entry{Macroeconomics (PhD)}{Fall 2023}
\role{Teaching Assistant to Prof. Seunghoon Na}{Purdue University}

\entry{Econometrics (Graduate)}{2019 -- 2021}
\role{TA to Profs. Kowsar Yousefi, Atiyeh Vahidmanesh, and Alireza Oryoie}{University of Tehran}

\entry{Game Theory (Graduate)}{Fall 2019}
\role{Teaching Assistant to Prof. Ghahreman Abdoli}{University of Tehran}

\entry{Econometrics in Transportation (Graduate)}{2020 -- 2021}
\role{Teaching Assistant to Prof. Amir Samimi}{Sharif University of Technology}

% ======================= AWARDS =======================
\section{Fellowships and Awards}
\begin{tight}
  \item Bilsland Dissertation Fellowship, Purdue University \hfill 2026--27
  \item Outstanding Teaching Award (ECON 252), Purdue University \hfill Fall 2025
  \item Distinguished Teaching Award (ECON 252), Purdue University \hfill Summer 2025
\end{tight}

% ======================= SERVICE =======================
\section{Service}
\begin{tight}
  \item Economics Senator, Business Doctoral Students Association, Purdue University \hfill % TODO: years
\end{tight}

% ======================= SKILLS =======================
\section{Programming and Data Skills}
Python, R, STATA, MATLAB, Claude Code\par\vspace{0.3em}
Structural modeling, numerical equilibrium computation, simulation-based methods.

% ======================= LANGUAGES =======================
\section{Languages}
Farsi (Native), English (Fluent)

% ======================= REFERENCES =======================
\section{References}

\begin{tabular}{@{}p{0.47\textwidth}p{0.47\textwidth}@{}}
\textbf{Cathy Zhang} (Chair) & \textbf{Mario J. Crucini}\\
Professor & Professor\\
Mitch Daniels School of Business & Mitch Daniels School of Business\\
Purdue University & Purdue University\\
\href{mailto:cmzhang@purdue.edu}{cmzhang@purdue.edu} & \href{mailto:mcrucini@purdue.edu}{mcrucini@purdue.edu}\\[0.8em]
\textbf{Chong Xiang} & \\
Professor & \\
Mitch Daniels School of Business & \\
Purdue University & \\
\href{mailto:cxiang@purdue.edu}{cxiang@purdue.edu} & \\
\end{tabular}

\vfill
\begin{center}\footnotesize Updated \monthyear\end{center}

\end{document}